% 分野別類題: オイラー方程式 解答例
% コンパイル: TEXINPUTS="../../../00_common:$TEXINPUTS" latexmk -xelatex answers.tex
\documentclass[a4paper,11pt]{article}
\usepackage{preamble}
\usepackage{shoakubox}

\begin{document}

\kamokumei{類題演習 解答例 --- オイラー方程式}

\daimon{【解法】オイラー方程式の解き方と導出}

\noindent
オイラー方程式（コーシー・オイラー方程式）とは、次の形の2階線形微分方程式である（$x>0$ とする）。
\[
x^{2}y'' + axy' + by = f(x)
\]
係数が $x$ のべきになっているため、$x = e^{t}$（すなわち $t = \ln x$）と置換すると定数係数の線形微分方程式に帰着できる。

\medskip
\noindent
\textbf{置換の導出。}
$t = \ln x$ とすると $\dfrac{dt}{dx} = \dfrac{1}{x}$ なので、合成関数の微分により
\[
\frac{dy}{dx} = \frac{dy}{dt}\cdot\frac{dt}{dx} = \frac{1}{x}\frac{dy}{dt}
\quad\Longrightarrow\quad
x\frac{dy}{dx} = \frac{dy}{dt}
\]
さらに2階微分を計算すると
\[
\frac{d^{2}y}{dx^{2}}
= \frac{d}{dx}\!\left(\frac{1}{x}\frac{dy}{dt}\right)
= -\frac{1}{x^{2}}\frac{dy}{dt} + \frac{1}{x}\cdot\frac{1}{x}\frac{d^{2}y}{dt^{2}}
= \frac{1}{x^{2}}\left(\frac{d^{2}y}{dt^{2}} - \frac{dy}{dt}\right)
\]
したがって
\[
x^{2}\frac{d^{2}y}{dx^{2}} = \frac{d^{2}y}{dt^{2}} - \frac{dy}{dt}
\]
これらを $D = \dfrac{d}{dt}$ と略記すれば、$x\dfrac{dy}{dx} = Dy$、$x^{2}\dfrac{d^{2}y}{dx^{2}} = D(D-1)y$ となる。よってオイラー方程式は
\[
D(D-1)y + aDy + by = f(e^{t})
\quad\Longrightarrow\quad
\frac{d^{2}y}{dt^{2}} + (a-1)\frac{dy}{dt} + by = f(e^{t})
\]
という $t$ に関する定数係数線形微分方程式になる。

\medskip
\noindent
\textbf{同次方程式の特性方程式。}
同次形 $x^{2}y''+axy'+by=0$ に対応する特性方程式は
\[
m(m-1) + am + b = 0
\]
であり、これは元の方程式に $y = x^{m}$ を直接代入しても得られる（$y'=mx^{m-1}$, $y''=m(m-1)x^{m-2}$ を代入すると $x^{m}\left[m(m-1)+am+b\right]=0$）。
\begin{itemize}
\item 異なる実根 $m_{1}, m_{2}$ のとき: $y = C_{1}x^{m_{1}} + C_{2}x^{m_{2}}$
\item 重根 $m$ のとき: $y = (C_{1} + C_{2}\ln x)\,x^{m}$
\item 共役複素根 $m = \alpha \pm i\beta$ のとき: $y = x^{\alpha}\left(C_{1}\cos(\beta\ln x) + C_{2}\sin(\beta\ln x)\right)$
\end{itemize}
（複素根の場合は $x^{i\beta} = e^{i\beta\ln x} = \cos(\beta\ln x) + i\sin(\beta\ln x)$ であることによる。）

非同次の場合は、$t$ に関する定数係数方程式に直したうえで未定係数法や定数変化法により特解 $y_{p}(t)$ を求め、$t=\ln x$ を戻して一般解 $y = y_{h} + y_{p}$ とすればよい。

\clearpage
\begin{mondaiboxS}{問1}
$x^{2}y'' - 2xy' + 2y = 0 \quad (x>0)$ の一般解を求めよ。
\end{mondaiboxS}
\kaitou
特性方程式は
\[
m(m-1) - 2m + 2 = 0
\quad\Longrightarrow\quad
m^{2} - 3m + 2 = 0
\quad\Longrightarrow\quad
(m-1)(m-2) = 0
\]
より $m = 1, 2$（異なる実根）。よって一般解は
\[
\boxed{\; y = C_{1}x + C_{2}x^{2} \;}
\qquad (C_{1}, C_{2} \text{ は任意定数})
\]
（検算: $y=x$ のとき $y'=1,\ y''=0$ で $x^{2}\cdot 0 - 2x\cdot 1 + 2x = 0$。$y=x^{2}$ のとき $y'=2x,\ y''=2$ で $x^{2}\cdot 2 - 2x\cdot 2x + 2x^{2} = 2x^{2}-4x^{2}+2x^{2}=0$。ともに成立。）

\clearpage
\begin{mondaiboxS}{問2}
$x^{2}y'' + xy' + y = 0 \quad (x>0)$ の一般解を求めよ。
\end{mondaiboxS}
\kaitou
特性方程式は
\[
m(m-1) + m + 1 = 0
\quad\Longrightarrow\quad
m^{2} + 1 = 0
\quad\Longrightarrow\quad
m = \pm i
\]
共役複素根 $\alpha = 0,\ \beta = 1$ であるから、一般解は
\[
\boxed{\; y = C_{1}\cos(\ln x) + C_{2}\sin(\ln x) \;}
\qquad (C_{1}, C_{2} \text{ は任意定数})
\]
（検算: $y=\cos(\ln x)$ のとき $y' = -\dfrac{\sin(\ln x)}{x}$, $y'' = \dfrac{-\cos(\ln x) + \sin(\ln x)}{x^{2}}$ より
$x^{2}y''+xy'+y = \left[-\cos(\ln x)+\sin(\ln x)\right] - \sin(\ln x) + \cos(\ln x) = 0$ となり成立。$\sin(\ln x)$ についても同様に成立することが確かめられる。）

\clearpage
\begin{mondaiboxS}{問3}
$x^{2}y'' - xy' + y = \ln x, \quad y(1) = 1,\ y'(1) = 0 \quad (x>0)$ を解け。
\end{mondaiboxS}
\kaitou
$t = \ln x$ とおくと、解法冒頭の変換により
\[
\frac{d^{2}y}{dt^{2}} - 2\frac{dy}{dt} + y = t
\]
となる（$a=-1,\ b=1$ より $a-1=-2$、右辺は $f(e^{t}) = \ln(e^{t}) = t$）。

まず斉次方程式の特性方程式は $m^{2}-2m+1=0$ より $(m-1)^{2}=0$、重根 $m=1$。よって
\[
y_{h} = (C_{1}+C_{2}t)e^{t}
\]
次に特解を $y_{p} = At+B$（$t=0,1$ は特性方程式の根でないので多項式のまま仮定できる）とおくと、$y_{p}'=A,\ y_{p}''=0$ であるから
\[
0 - 2A + (At+B) = t
\quad\Longrightarrow\quad
At + (B-2A) = t
\]
係数比較により $A=1,\ B=2A=2$。よって $y_{p} = t+2$。

したがって $t$ に関する一般解は $y = (C_{1}+C_{2}t)e^{t} + t + 2$。$t=\ln x,\ e^{t}=x$ を代入して $x$ に戻すと
\[
y = C_{1}x + C_{2}x\ln x + \ln x + 2
\]
初期条件を適用する。$y' = C_{1} + C_{2}(\ln x + 1) + \dfrac{1}{x}$ である。

$x=1$（$\ln 1 = 0$）で
\[
y(1) = C_{1} + 2 = 1 \implies C_{1} = -1
\]
\[
y'(1) = C_{1} + C_{2} + 1 = 0 \implies -1 + C_{2} + 1 = 0 \implies C_{2} = 0
\]
よって
\[
\boxed{\; y = 2 - x + \ln x \;}
\]
（検算: $y'=-1+\dfrac{1}{x}$, $y''=-\dfrac{1}{x^{2}}$ より
$x^{2}y'' - xy' + y = -1 - x\left(-1+\dfrac{1}{x}\right) + (2-x+\ln x) = -1 + x - 1 + 2 - x + \ln x = \ln x$ で方程式を満たす。
また $y(1)=2-1+0=1$、$y'(1)=-1+1=0$ で初期条件も満たす。）

\end{document}
